%%%%%%%%%%%%%%%%%%%%%%%%%%%%%%%%%%%%%%%%%%%%%%%%%%%%%%%%%%%%%%%%%%%%%%%%
\section{The conjectured meander limit}\label{sec:limit}
%%%%%%%%%%%%%%%%%%%%%%%%%%%%%%%%%%%%%%%%%%%%%%%%%%%%%%%%%%%%%%%%%%%%%%%%

The APPS model concerns determinants, not curves. For meanders themselves, Borga, Gwynne and Sun
proposed a precise description of a large random meander seen from far away~\cite{BGS}. This
section explains it; nothing in it is new. Close the river into a loop through infinity. River and
road are then two loops on the sphere crossing $2n$ times, and their union is a planar map whose
$2n$ vertices have degree four and are each visited once by both curves. A uniform meander of order
$n$ is thus a random planar map with two curves through every vertex, and each curve orders the
vertices.

\subsection{The surface and the curves}

A \emph{Liouville quantum gravity sphere} with parameter $\gamma\in(0,2)$ is, for our purposes, a
random measure of total mass one on the sphere, of the form $e^{\gamma h}d^2z$ for a variant $h$ of the Gaussian free
field, defined as a Gaussian multiplicative chaos as in Section~\ref{sec:upper}. The variant has
logarithmic singularities at marked points and a prescribed behavior at infinity, and its law is
weighted by a power of its total mass before the mass is normalized to one~\cite{DMS},
\cite[Section~2]{BGS}. We write $\nu$ for the area. Such spheres are the expected limits of random
planar maps. For uniform planar maps ($c=0$, $\gamma^2=8/3$) this is proved in the sense of metric
spaces: the limit of the graph distance, suitably rescaled, is the \emph{Brownian map}, a random
compact metric space~\cite{LeGall,Miermont}. For meanders the parameter is the value that~\eqref{eq:Q} assigns to
$c=-4$, namely $\gamma^2=(17-\sqrt{145})/3\approx1.6528$.

\emph{Schramm--Loewner evolutions} SLE$_\kappa$ are a one-parameter family of random fractal
curves in the plane, characterized by invariance under conformal maps and a Markov property: given
an initial piece, the rest of the curve is again an SLE$_\kappa$ in the remaining domain. The
parameter $\kappa>0$ measures how much the curve wiggles. For $\kappa\ge8$ the curve passes through
every point. SLE$_8$ is the scaling limit of the curve that runs around a uniformly
random spanning tree~\cite{LSWust}. Its whole-plane version starts and ends at infinity and visits
every point of the sphere; a fixed point is visited exactly once, almost surely. Given the surface,
a curve is \emph{traced at unit area speed} if it is parametrized by $[0,1]$ so that the region
covered during $[0,t]$ has area $t$. The two curves of the limit object are chosen independently of
each other and of the field before this parametrization; afterwards they share the area clock.

\begin{conjecture}[Borga, Gwynne and Sun {\cite[Conjecture~1.4]{BGS}}]\label{conj:bgs}
Large uniformly random meanders, with the uniform measure on vertices and a suitable rescaling of
distances, converge in law to a unit-area Liouville quantum gravity sphere with
$\gamma^2=(17-\sqrt{145})/3$, decorated by two independent whole-plane space-filling SLE$_8$
curves traced at unit area speed.
\end{conjecture}

A weaker form concerns only the orders of the crossings. A \emph{permuton} is a
probability measure on the unit square $[0,1]^2$ whose two marginals are uniform. A permutation
$\omega$ of $\{1,\ldots,2n\}$ gives one: spread mass $1/(2n)$ uniformly over each small square
$[\frac{i-1}{2n},\frac{i}{2n}]\times[\frac{\omega(i)-1}{2n},\frac{\omega(i)}{2n}]$. Permutons are the
natural limits of large permutations. For a meander, number the crossings $1,\ldots,2n$ from left
to right along the river, and let $\omega(i)$ be the position of crossing $i$ along the road,
traversed counterclockwise starting from crossing $1$, so $\omega(1)=1$~\cite[Definition~1.6]{BGS}.
In the limit object both curves start and end at the marked point $\infty$, which plays the role of
crossing $1$. Another starting crossing or direction would only rotate or reflect the cyclic road
order, which the events of Section~\ref{sec:nearsplit} do not see. In the limit object, the analogous
measure is the \emph{meandric permuton}
\[
\Pi(\mathcal A\times\mathcal B)=\nu\bigl(\eta_1(\mathcal A)\cap\eta_2(\mathcal B)\bigr),\qquad \mathcal A,\mathcal B\subset[0,1],
\]
where $\eta_1$ is the river curve and $\eta_2$ the road curve: $\Pi(\mathcal A\times\mathcal B)$ is the area of the
set of points that the river visits during the times $\mathcal A$ and the road during the times $\mathcal B$. Its
marginals are uniform because each curve is traced at unit area speed.

\begin{conjecture}[Borga, Gwynne and Sun {\cite[Conjecture~1.8]{BGS}}]\label{conj:permuton}
The permuton of a uniformly random meander of order $n$ converges in law, as $n\to\infty$, to the
meandric permuton $\Pi$.
\end{conjecture}

Both conjectures are open. BGS compute properties of the candidate limit, such as the expected
frequencies of small patterns in the permutation, but neither convergence nor the exponent
follows.

\emph{Why the conjecture is hard.} In the theory of Duplantier, Miller and Sheffield (``mating
of trees''), a $\gamma$-LQG surface and an independent
space-filling SLE$_\kappa$ with $\kappa=16/\gamma^2$ together have a simple description by a pair
of correlated Brownian motions. Call such a curve \emph{matched} to the surface. For one spanning
tree, $\gamma^2=2$ and $\kappa=8$ match. Mullin's bijection encodes a map with a spanning tree by two
walks, one tracing the tree and one its dual, and these walks converge to the pair of Brownian
motions~\cite{Sheffield,DMS}. For meanders each curve alone would call for $\gamma^2=2$, while the
pair calls for $\gamma^2\approx1.653$, whose matched curve has $\kappa=4\xstar\approx9.68$. No current
method handles curves that are not matched to the surface. The
next two sections therefore prove results about the limit object itself, which is well defined,
without any claim that meanders converge to it.
